
A finalidade deste capítulo é delimitar o que se pretende construir ao longo das duas fases do Trabalho de Conclusão de Curso e, em seguida, apresentar a cadeia temporal de atividades prevista na pré-proposta aprovada. A descrição parte da plataforma vista como produto educacional, atravessa as macroatividades técnicas necessárias à sua materialização e culmina na exposição do cronograma e do regime de trabalho em dupla adotado.

O artefato que orienta este trabalho não é apenas um compilador ou um editor de código, mas sim uma plataforma educacional web cuja unidade de valor é a possibilidade de o estudante personalizar a própria linguagem que utilizará para aprender programação. Em termos concretos, oferece-se ao usuário a possibilidade de redefinir o vocabulário das palavras-chave, escolher os delimitadores de bloco que melhor representem sua noção de escopo, substituir operadores lógicos e relacionais por palavras de uso natural, reescrever os literais booleanos no idioma de sua preferência e ajustar as regras de finalização de instrução. Cada uma dessas dimensões responde, em alguma medida, a uma fonte conhecida de atrito sintático no contato inicial com linguagens consolidadas no mercado, sendo a sua reunião sob uma mesma interface o que sustenta o caráter inovador da proposta. A Tabela~\ref{tab:elementos-configuraveis} detalha cada um desses elementos, confrontando a forma canônica herdada da família C com um exemplo de personalização e com as restrições verificadas pela plataforma no momento da configuração.

\begin{table}[htbp]
    \centering
    \caption{Elementos da linguagem passíveis de redefinição pelo usuário.}
    \label{tab:elementos-configuraveis}
    \footnotesize
    \setlength{\tabcolsep}{4pt}
    \renewcommand{\arraystretch}{1.25}
    \begin{tabular}{|>{\raggedright\arraybackslash}p{2.5cm}|>{\raggedright\arraybackslash}p{2.6cm}|>{\raggedright\arraybackslash}p{2.9cm}|>{\raggedright\arraybackslash}p{5.2cm}|}
        \hline
        \textbf{Elemento} & \textbf{Forma canônica} & \textbf{Exemplo de personalização} & \textbf{Restrições verificadas} \\
        \hline
        Palavras-chave de tipo & \texttt{int}, \texttt{float}, \texttt{bool}, \texttt{string} & \texttt{inteiro}, \texttt{real}, \texttt{logico}, \texttt{texto} & Devem ter forma de identificador e não colidir entre si nem com os demais elementos configurados \\
        \hline
        Controle de fluxo & \texttt{if}, \texttt{else}, \texttt{for}, \texttt{while}, \texttt{break}, \texttt{continue}, \texttt{return} & \texttt{se}, \texttt{senao}, \texttt{para}, \texttt{enquanto}, \texttt{pare}, \texttt{siga}, \texttt{retorne} & Mesmas restrições das palavras-chave de tipo \\
        \hline
        Entrada e saída & \texttt{print}, \texttt{scan} & \texttt{escreva}, \texttt{leia} & Mesmas restrições das palavras-chave de tipo \\
        \hline
        Operadores lógicos & \texttt{\&\&}, \texttt{||}, \texttt{!} & \texttt{e}, \texttt{ou}, \texttt{nao} & Forma de identificador; sem duplicatas; não podem reusar palavra reservada nem colidir com delimitadores de bloco \\
        \hline
        Operadores relacionais & \texttt{==}, \texttt{!=}, \texttt{>}, \texttt{>=}, \texttt{<}, \texttt{<=} & \texttt{igual}, \texttt{diferente}, \texttt{maior}, \texttt{maiorIgual}, \texttt{menor}, \texttt{menorIgual} & Mesmas restrições dos operadores lógicos \\
        \hline
        Literais booleanos & \texttt{true}, \texttt{false} & \texttt{verdadeiro}, \texttt{falso} & Forma de identificador; sem duplicatas; não podem conflitar com operadores em forma de palavra nem com delimitadores \\
        \hline
        Delimitadores de bloco & \texttt{\{} e \texttt{\}} & \texttt{inicio} e \texttt{fim}, ou indentação & Ambos preenchidos, com forma de identificador, distintos entre si e sem reusar palavra reservada \\
        \hline
        Terminador de instrução & \texttt{;} & \texttt{ponto} ou ausência de terminador & Não pode ser vazio, conter espaços, reusar o ponto e vírgula ou caracteres de operadores fixos, nem colidir com palavras já definidas na linguagem \\
        \hline
    \end{tabular}
    \legend{Fonte: elaborado pelo autor, a partir das funções de validação do módulo \texttt{config.ts} do pacote \texttt{compiler}.}
\end{table}


O conjunto completo de funcionalidades previstas para a plataforma, com os respectivos atores, é apresentado no diagrama de casos de uso da Figura~\ref{fig:casos-de-uso}. O diagrama distingue quatro atores: o visitante, que apenas se cadastra e se autentica; o usuário autenticado, que concentra as funcionalidades comuns de personalização da linguagem e de uso do ambiente de desenvolvimento; e os papéis de aluno e de professor, que dele herdam e acrescentam, respectivamente, as funcionalidades de participação em turmas e de gestão pedagógica.

\begin{landscape}
\begin{figure}[p]
    \centering
    \caption{Diagrama de casos de uso da plataforma.}
    \label{fig:casos-de-uso}
    \resizebox{\linewidth}{!}{%
    \begin{tikzpicture}[x=1cm,y=1cm]

    %% ---------- fronteira do sistema ----------
    \draw[thick] (3.0,5.9) rectangle (23.2,-7.2);
    \node[font=\bfseries, anchor=north] at (13.1,5.7) {Plataforma educacional web};

    %% ---------- atores ----------
    \pic at (0,4.6) {ator};  \node[font=\small, align=center] at (0,2.9) {Visitante};
    \pic at (0,0.9) {ator};  \node[font=\small, align=center] at (0,-1.1) {Usuário\\autenticado};
    \pic at (0,-4.4) {ator}; \node[font=\small, align=center] at (0,-6.1) {Aluno};
    \pic at (26.0,-0.6) {ator};  \node[font=\small, align=center] at (26.0,-2.4) {Professor};

    %% generalizacao: Aluno e Professor sao Usuarios autenticados.
    %% A do Professor e roteada abaixo da fronteira para nao cruzar os casos de uso.
    \draw[heranca] (0,-4.2) -- (0,-2.0);
    \draw[heranca] (26.0,-1.9) -- (26.0,-8.3) -- (1.5,-8.3) -- (1.5,0.3) -- (0.45,0.3);

    %% ---------- coluna A: acesso, conta e linguagem ----------
    \node[uc] (ucCad)  at (7.2, 4.5)  {Cadastrar-se};
    \node[uc] (ucAut)  at (7.2, 3.0)  {Autenticar-se};
    \node[uc] (ucPer)  at (7.2, 1.5)  {Gerenciar perfil};
    \node[uc] (ucLing) at (7.2, 0.0)  {Criar e gerenciar\\linguagens personalizadas};
    \node[uc] (ucCom)  at (7.2,-1.5)  {Compartilhar e importar\\linguagens da comunidade};
    \node[uc] (ucAtiv) at (7.2,-3.0)  {Ativar linguagem\\no editor};
    \node[uc] (ucEdit) at (7.2,-4.5)  {Editar código na IDE};
    \node[uc] (ucArq)  at (7.2,-6.0)  {Gerenciar arquivos\\do projeto};

    %% ---------- coluna B: ambiente de desenvolvimento e aluno ----------
    \node[uc] (ucExec) at (13.6, 4.5) {Executar programa\\no terminal};
    \node[uc] (ucTok)  at (13.6, 3.0) {Inspecionar \textit{tokens} e\\código intermediário};
    \node[uc] (ucDep)  at (13.6, 1.5) {Depurar execução\\passo a passo};
    \node[uc] (ucIng)  at (13.6, 0.0) {Ingressar em turma};
    \node[uc] (ucCons) at (13.6,-1.5) {Consultar listas\\e exercícios};
    \node[uc] (ucSub)  at (13.6,-3.0) {Submeter solução\\de exercício};
    \node[uc] (ucVer)  at (13.6,-4.5) {Consultar veredito\\e histórico};

    %% ---------- coluna C: gestao pedagogica ----------
    \node[uc] (ucTur)  at (20.0, 4.5) {Gerenciar turmas};
    \node[uc] (ucMem)  at (20.0, 3.0) {Gerenciar membros\\da turma};
    \node[uc] (ucEx)   at (20.0, 1.5) {Gerenciar exercícios\\e casos de teste};
    \node[uc] (ucVinc) at (20.0, 0.0) {Vincular linguagem\\a exercício};
    \node[uc] (ucLis)  at (20.0,-1.5) {Gerenciar listas\\de exercícios};
    \node[uc] (ucPub)  at (20.0,-3.0) {Publicar lista\\para a turma};
    \node[uc] (ucAva)  at (20.0,-4.5) {Acompanhar e avaliar\\submissões};

    %% ---------- associacoes (camada de fundo: passam atras das elipses) ----------
    \begin{scope}[on background layer]
      \foreach \d in {ucCad,ucAut}
        \draw[liga] (0.35,4.2) -- (\d);
      \foreach \d in {ucPer,ucLing,ucCom,ucAtiv,ucEdit,ucArq,ucExec,ucTok,ucDep}
        \draw[liga] (0.35,0.5) -- (\d);
      \foreach \d in {ucIng,ucCons,ucSub,ucVer}
        \draw[liga] (0.35,-4.8) -- (\d);
      \foreach \d in {ucTur,ucMem,ucEx,ucVinc,ucLis,ucPub,ucAva}
        \draw[liga] (25.3,-0.9) -- (\d);
    \end{scope}

    \end{tikzpicture}}
    \legend{Fonte: elaborado pelo autor.}
\end{figure}
\end{landscape}

A partir desses casos de uso, consolidaram-se os requisitos que orientaram a construção da plataforma. A Tabela~\ref{tab:requisitos-funcionais} relaciona os requisitos funcionais, isto é, o que o sistema deve fazer; a Tabela~\ref{tab:requisitos-nao-funcionais}, os requisitos não funcionais, que estabelecem as qualidades esperadas da solução e os critérios sob os quais ela será verificada.

\begin{table}[htbp]
    \centering
    \caption{Requisitos funcionais da plataforma.}
    \label{tab:requisitos-funcionais}
    \footnotesize
    \setlength{\tabcolsep}{4pt}
    \renewcommand{\arraystretch}{1.15}
    \begin{tabular}{|c|>{\raggedright\arraybackslash}p{12.3cm}|}
        \hline
        \textbf{ID} & \textbf{Descrição} \\
        \hline
        RF01 & Permitir o cadastro e a autenticação de usuários, distinguindo os papéis de aluno e de professor \\
        \hline
        RF02 & Permitir ao usuário autenticado consultar e editar os dados de seu perfil \\
        \hline
        RF03 & Oferecer um assistente guiado, organizado em etapas independentes, para a criação de uma linguagem personalizada \\
        \hline
        RF04 & Permitir a redefinição dos elementos da linguagem relacionados na Tabela~\ref{tab:elementos-configuraveis} \\
        \hline
        RF05 & Validar preventivamente cada configuração, impedindo colisões entre palavras-chave, operadores, literais, delimitadores e terminador de instrução \\
        \hline
        RF06 & Exibir, em tempo real, a pré-visualização de um programa exemplo sob a configuração corrente \\
        \hline
        RF07 & Persistir localmente a customização ativa e remotamente as linguagens nomeadas pelo usuário \\
        \hline
        RF08 & Permitir publicar uma linguagem no catálogo comunitário e importar linguagens publicadas por terceiros \\
        \hline
        RF09 & Oferecer edição de código com realce de sintaxe e sugestão automática ajustados à linguagem ativa \\
        \hline
        RF10 & Manter um sistema de arquivos virtual que permita criar, renomear, remover e alternar entre arquivos \\
        \hline
        RF11 & Executar o programa produzido no editor, com entrada e saída conectadas ao terminal integrado \\
        \hline
        RF12 & Exibir ao usuário a tabela de \textit{tokens} reconhecidos e o código intermediário gerado \\
        \hline
        RF13 & Oferecer execução passo a passo, destacando no editor a instrução em processamento \\
        \hline
        RF14 & Relatar erros léxicos, sintáticos, semânticos e de execução no vocabulário da linguagem ativa e no idioma selecionado \\
        \hline
        RF15 & Permitir ao professor criar e gerir turmas, e ao aluno ingressar em uma turma \\
        \hline
        RF16 & Permitir ao professor criar e gerir exercícios, com casos de teste de entrada e saída esperada \\
        \hline
        RF17 & Permitir vincular uma configuração de linguagem a um exercício, restringindo o vocabulário disponível ao aluno \\
        \hline
        RF18 & Permitir organizar exercícios em listas e publicá-las para uma turma \\
        \hline
        RF19 & Avaliar automaticamente a submissão do aluno contra os casos de teste do exercício \\
        \hline
        RF20 & Permitir ao professor acompanhar as submissões da turma e atribuir nota \\
        \hline
    \end{tabular}
    \legend{Fonte: elaborado pelo autor.}
\end{table}

\begin{table}[htbp]
    \centering
    \caption{Requisitos não funcionais da plataforma.}
    \label{tab:requisitos-nao-funcionais}
    \footnotesize
    \setlength{\tabcolsep}{4pt}
    \renewcommand{\arraystretch}{1.15}
    \begin{tabular}{|c|>{\raggedright\arraybackslash}p{12.3cm}|}
        \hline
        \textbf{ID} & \textbf{Descrição} \\
        \hline
        RNF01 & \textbf{Acessibilidade:} os componentes de interface devem observar as diretrizes WAI-ARIA \cite{w3c2023aria}, com navegação por teclado, leitura por tecnologias assistivas e gerenciamento de foco \\
        \hline
        RNF02 & \textbf{Responsividade:} as composições de tela devem reorganizar-se conforme a faixa de resolução do dispositivo \\
        \hline
        RNF03 & \textbf{Internacionalização:} toda cadeia visível ao usuário deve estar externalizada, com suporte aos idiomas \textit{pt-BR}, \textit{pt-PT}, \textit{es} e \textit{en} \\
        \hline
        RNF04 & \textbf{Autonomia de execução:} edição, análise léxica e sintática, geração de código intermediário e interpretação devem permanecer funcionais na ausência de conexão com o serviço de persistência \\
        \hline
        RNF05 & \textbf{Latência:} a execução do núcleo de tradução deve ocorrer sem ciclo de requisição e resposta a servidor remoto \\
        \hline
        RNF06 & \textbf{Consistência visual:} os temas claro e escuro devem manter paridade de contraste e legibilidade, inclusive no editor \\
        \hline
        RNF07 & \textbf{Testabilidade:} as construções críticas devem ser cobertas por testes automatizados, executados localmente e no \textit{pipeline} de integração contínua \\
        \hline
        RNF08 & \textbf{Segregação de acesso:} as rotas e operações restritas devem verificar o papel do usuário antes de permitir o acesso \\
        \hline
    \end{tabular}
    \legend{Fonte: elaborado pelo autor.}
\end{table}


Conforme definido na pré-proposta, a metodologia foi organizada em seis macroetapas complementares. Essas etapas não devem ser entendidas como blocos completamente isolados, pois o desenvolvimento de uma plataforma interativa exige revisões sucessivas entre interface, interpretador e experiência de uso. Ainda assim, elas delimitam com clareza o percurso técnico e acadêmico do projeto:

\begin{itemize}
    \item \textbf{Revisão bibliográfica:} levantamento das bases conceituais sobre interpretadores, compiladores e ferramentas educacionais relacionadas, com o objetivo de sustentar tanto as decisões de projeto quanto a posterior análise comparativa;
    \item \textbf{Definição de requisitos e arquitetura:} consolidação dos elementos configuráveis da linguagem, da modelagem das camadas da plataforma e dos contratos de comunicação entre o núcleo do interpretador e a IDE;
    \item \textbf{Desenvolvimento da plataforma:} implementação da aplicação web, contemplando os fluxos de edição, navegação, personalização, autenticação e uso educacional da IDE;
    \item \textbf{Implementação da análise dinâmica léxica e sintática:} integração entre o código escrito na IDE e as fases de reconhecimento de \textit{tokens} e validação sintática com suporte às regras configuradas pelo usuário;
    \item \textbf{Geração e interpretação do código intermediário:} disponibilização, no ambiente da IDE, da execução do código intermediário e do retorno de avisos, erros e resultados ao usuário;
    \item \textbf{Integração e testes entre plataforma e interpretador:} validação dos fluxos principais, com atenção à consistência da experiência, à acessibilidade, à responsividade e à preparação para os estudos com usuários.
\end{itemize}

A distribuição temporal dessas atividades ao longo dos doze meses que compreendem o TCC~I e o TCC~II é apresentada na Tabela~\ref{tab:cronograma}, conforme o cronograma registrado na pré-proposta. Cada coluna numerada corresponde a um mês corrido das duas etapas e cada célula marcada com ``X'' sinaliza a previsão de execução da respectiva atividade no período indicado.

\begin{table}[H]
    \centering
    \caption{Cronograma de desenvolvimento do trabalho de conclusão de curso.}
    \label{tab:cronograma}
    \footnotesize
    \setlength{\tabcolsep}{4pt}
    \renewcommand{\arraystretch}{1.2}
    \resizebox{\textwidth}{!}{%
    \begin{tabular}{|p{6.5cm}|c|c|c|c|c|c|c|c|c|c|c|c|}
        \hline
        \multirow{2}{*}{\textbf{Atividade}} & \multicolumn{6}{c|}{\textbf{TCC I}} & \multicolumn{6}{c|}{\textbf{TCC II}} \\
        \cline{2-13}
         & \textbf{01} & \textbf{02} & \textbf{03} & \textbf{04} & \textbf{05} & \textbf{06} & \textbf{07} & \textbf{08} & \textbf{09} & \textbf{10} & \textbf{11} & \textbf{12} \\
        \hline
        Revisão Bibliográfica & X & X & X & X & & & & & & & & \\
        \hline
        Escrita da Introdução e Revisão de Literatura & & X & X & X & & & & & & & & \\
        \hline
        Predefinição de escopo e Especificação dos comandos dinâmicos & & & X & X & & & & & & & & \\
        \hline
        Escrita da Metodologia & & & X & & X & X & & & & & & \\
        \hline
        Arquitetura e Modelagem da Plataforma & & & X & X & & & & & & & & \\
        \hline
        Customização dinâmica de \textit{tokens} e sintaxe & & & & X & X & & & & & & & \\
        \hline
        Execução da Análise Léxica e Sintática com o código da IDE & & & & & X & X & & & & & & \\
        \hline
        Execução do código intermediário no Terminal da IDE com \textit{feedback} (avisos e erros) & & & & & & X & X & X & & & & \\
        \hline
        Finalização de fluxos, acessibilidade e responsividade & & & & & & X & X & X & & & & \\
        \hline
        Testes, Validação com Usuários e Coleta de \textit{Feedback} & & & & & & & & & X & X & X & \\
        \hline
        Escrita do TCC (Resultados, Discussão e Conclusão) & & & & & & & & & X & X & X & X \\
        \hline
    \end{tabular}%
    }
    \legend{Fonte: elaborado pelo autor.}
\end{table}

Convém retomar, neste ponto, o regime de trabalho em dupla já apresentado na introdução (Capítulo~\ref{cap:introducao_modelo}), dada sua influência direta sobre o escopo do projeto e sobre os capítulos subsequentes. A pré-proposta aprovada para este Trabalho de Conclusão de Curso prevê a \textbf{participação conjunta em todas as etapas} do desenvolvimento, mas com funções complementares. Em respeito a esse compromisso, as decisões de arquitetura, os ciclos de revisão de código e as etapas de integração foram conduzidos de modo colaborativo pelos dois integrantes da dupla. Cabe ressaltar que, por se tratar de uma equipe reduzida, não houve especialização de papéis: ambos os integrantes acumularam as atribuições de levantamento de requisitos, projeto, implementação, teste e revisão. Adotaram-se, de forma pragmática, as práticas técnicas do \textit{Extreme Programming} discutidas no referencial teórico, sem ciclos de duração fixa: partiu-se de um conjunto mínimo de requisitos e o sistema avançou por acréscimos funcionais, incorporando as melhorias identificadas durante a própria implementação \cite{beck2004,wildt2015}.

A partir dessa base comum, a divisão complementar de eixos de aprofundamento manifesta-se nos capítulos de redação acadêmica e nas evoluções incrementais subsequentes. No presente documento, o autor concentra a investigação na plataforma e na customização de comandos: arquitetura da aplicação web, concepção do assistente de personalização, experiência de uso da IDE, execução do código no terminal integrado e procedimentos de coleta de \textit{feedback} junto aos estudantes. O detalhamento aprofundado da evolução do núcleo do interpretador, incluindo a geração de código intermediário e sua integração técnica ao restante da plataforma, é desenvolvido em paralelo pelo estudante Victor de Souza Vilela da Silva. Essa divisão de eixos não fragmenta o produto: o que se distribui é a profundidade analítica de cada autor sobre partes específicas do mesmo artefato, conduzido pelos dois em todas as etapas. Tal arranjo justifica-se pela amplitude do escopo, que, individualmente, tenderia a resultar em um interpretador simplificado ou em uma plataforma incompleta, restrita a um terminal.

Concluídas as etapas cobertas por este documento, dedicadas à fundamentação teórica, à modelagem da plataforma, à customização dinâmica de \textit{tokens} e sintaxe, à execução das análises léxica e sintática a partir da IDE e à execução do código intermediário no terminal integrado, a fase subsequente direcionará seus esforços à consolidação da plataforma e à definição do protocolo de avaliação: os instrumentos de coleta, as métricas e os critérios que permitirão, posteriormente, aferir em que medida a personalização da linguagem contribui para a aprendizagem. A aplicação desse protocolo junto a estudantes, com a consequente coleta de evidências empíricas em sala de aula, não integra o escopo deste trabalho e é registrada como encaminhamento futuro. O próximo capítulo, dedicado à modelagem, descreve o estado em que o sistema se encontra ao término do TCC~I.
